\section{Uniform numerical parameters and polynomial bond bounds}
\label{sec:uniform_physical_soft_calibration}

Corollary~\ref{cor:calibrated_uniform_physical_soft_family} chooses the
parameters of Theorem~\ref{thm:uniform_physical_soft_improving_family}
before the chain length. It retains that theorem's actual family and
bounds its permitted member bond by one natural polynomial. The
parameter dependence is a corrected auxiliary to
\cite[Proposition~5.7 and Theorem~5.8]{Arad2016RigorousRG}, with the
qualifications in
\cite{gap:alvv16_soft_regional_residuals,gap:alvv16_cluster_approximation}.

\begin{lemma}[Uniform numerical soft parameters]
    \label{lem:uniform_soft_parameters}
    \lean{TNLean.MPS.AreaLaw.ALVV.exists_uniform_soft_parameters}
    \leanok
    Fix $q\ge2$ and $g>0$. There exist natural numbers
    $u\ge2$, $C_{\rm pol}\ge1$ and $p_{\rm pol}$ before every volume
    $N$, with the following properties. Define the quantities in
    (\ref{eq:calibrated_soft_parameters}), with $g$ in place of
    $\gamma$. Then $\ell\ge1$, $m\ge2$, $k\ge1$, $D\ge2$, $M\ge34$,
    $0<\tau,\xi\le1$, $t,T\ge156+2g$ and $T\ge g/40+24$.
    Both residual inequalities in~(\ref{eq:uniform_soft_residual_budgets})
    hold with $T_0=T_1=T$, as do the sampling inequalities
    (\ref{eq:calibrated_soft_sampling}) and, for every $N$, the member-bond
    bound (\ref{eq:calibrated_soft_member_bond}). These choices require
    neither an interaction Hamiltonian nor a sampled family.
\end{lemma}
\begin{proof}\leanok
    The uniform numerical calibration with exponents $(24,10,10)$
    chooses $u$ before the volume and interactions. The definitions give
    $D\ge2$ and $M\ge34$. Thus $D^{20}\ge10000$ and
    $M^{10}\ge100000$, so
    \[
      \tau\le\frac1{10000D^4},\qquad
      D^{12}\Delta\le\frac1{M^{10}}\le\frac1{100000}.
    \]
    These give the remaining sampling inequalities. The geometric-order
    polynomial estimate at base $2(e^{1/3}-1)$ and coefficient budget
    $t2^{m+1}$ gives natural $C_{\rm pol}$ and $p_{\rm pol}$ before $N$.
    Its coefficient multiplicity is $m+1$ and its filter degree is $k$.
    Since $B_N\ge1$, this is exactly the displayed member-bond bound.
\end{proof}

\begin{corollary}[A calibrated common physical soft family]
    \label{cor:calibrated_uniform_physical_soft_family}
    \lean{TNLean.MPS.AreaLaw.ALVV.exists_calibrated_uniform_physical_soft_improving_family}
    \leanok
    Fix $q\ge2$ and $\gamma>0$. There are natural numbers
    $u\ge2$, $C_{\rm pol}\ge1$ and $p_{\rm pol}$, before every chain
    length $N$, with the following properties. Put
    \begin{align}
      \ell&=u^8, & m&=u^{10}, & k&=u^{16}, &
      t&=3m, & T&=u^9,\notag\\
      D&=(mk+1)(qk)^{36\ell+4\lfloor k/\ell\rfloor}, &
      M&=3t(1+\log m)+4\ell+34,\notag\\
      \tau&=\frac1{D^{24}M^{10}}, &
      \xi&=\min\left\{1,\frac{\tau\gamma^2}{20000(M+\gamma+8)}\right\},\notag\\
      U&=M+\xi+4+\gamma, &
      \Delta&=4\exp\left(-4k\sqrt{\frac{\gamma}{2U}}\right).
      \label{eq:calibrated_soft_parameters}
    \end{align}
    The scale satisfies $t\ge156+2\gamma$, and the common residual
    cutoff satisfies $T\ge156+2\gamma$ and $T\ge\gamma/40+24$.
    Both inequalities in~(\ref{eq:uniform_soft_residual_budgets}) hold
    with $T_0=T_1=T$. In addition, $D\ge2$, $M\ge34$,
    $0<\tau,\xi\le1$, and
    \begin{align}
      \Delta&\le\frac1{D^{12}M^{10}},\notag\\
      \frac{3\tau}{200}
        &\le\frac1{100}\left(\frac{3}{100D^2}\right)^2, &
      \frac{3\tau}{200}&\le\frac1{200}, &
      D^{12}\Delta&\le\frac1{100000}.
      \label{eq:calibrated_soft_sampling}
    \end{align}

    For every $N$, let $s_N$ and $B_N=2s_N^2q^{s_N}(m+1)$ be exactly
    the cluster order and native bond in
    (\ref{eq:uniform_soft_native_family}). Then
    \begin{align}
      (\max\{1,B_N\})^{2k}
        \le C_{\rm pol}(N+1)^{p_{\rm pol}}.
      \label{eq:calibrated_soft_member_bond}
    \end{align}
    For every positive-contraction nearest-neighbour Hamiltonian and
    natural two-window geometry of
    Theorem~\ref{thm:uniform_physical_soft_improving_family}, all that
    theorem's conclusions hold with these same numerical choices.
    Thus one actual original minimum tuple, native core, polynomial
    closing coefficients and common left family precede every allowed
    regional estimate. Each member has its original allowed bond
    $(\max\{1,B_N\})^{2k}$, bounded by
    (\ref{eq:calibrated_soft_member_bond}). The approximants, their own
    closed spectral clusters, actual Chebyshev operators and all ordered
    word powers are unchanged.
\end{corollary}
\begin{proof}\leanok
    \uses{lem:uniform_soft_parameters,thm:uniform_physical_soft_improving_family}
    Lemma~\ref{lem:uniform_soft_parameters} supplies the common numerical
    choices and polynomial member-bond bound.
    Apply Theorem~\ref{thm:uniform_physical_soft_improving_family}
    once to the supplied original interactions with the derived
    numerical budgets. Retain the same family and all its conclusions together
    with~(\ref{eq:calibrated_soft_member_bond}); no other approximant,
    target or family is selected.
\end{proof}

Here the soft order $m$ and filter degree $k$ are fixed before $N$;
the cluster order $s_N$ varies according to its logarithmic formula. Thus
(\ref{eq:calibrated_soft_member_bond}) is a polynomial bound in the
volume. No polynomial assertion is made when the soft order or filter degree
is allowed to vary independently with $N$. Sharp inverse-gap dependence,
an adaptive generating algorithm and running time are separate conclusions.
